\section{Discusión}\label{sec:discussion}

\subsection{Mecanismos}

La evidencia describe un salario mínimo que movió trabajadores y no salarios. La reasignación es mi resultado más robusto: la formalidad cae y el trabajo independiente aumenta, los desplazamientos crecen con la incidencia (\emph{bite}) regional del piso y los confirma un diseño de exposición continua con inferencia por aleatorización. El mecanismo natural es el clásico: un piso que se vuelve vinculante en el sector cubierto induce una reasignación hacia el sector no cubierto, el canal de desplazamiento que \citet{jaramillo_lopez_2006} documentan para el Perú y \citet{ham_2018_honduras} para Honduras. Su incidencia aparece exactamente donde ese mecanismo lo predice: entre los jóvenes, entre las mujeres y en los departamentos de alta incidencia, es decir, en los trabajadores y lugares donde el nuevo piso cortó más hondo en la distribución salarial. Las transiciones de panel de la Sección~\ref{sec:transitions} identifican el flujo detrás del cambio en el stock y precisan el mecanismo en un aspecto importante: la retención de los trabajadores formales en su puesto no cayó, mientras que la entrada a empleos formales desde la informalidad se redujo casi a la mitad, de modo que la reasignación operó a través de la creación de vínculos formales y no de su destrucción. El piso no expulsó a quienes ya estaban empleados; le cerró la puerta a nuevos vínculos formales para los trabajadores de baja educación. Es el mismo ajuste por el margen de las contrataciones que \citet{dustmann_2022} documentan para Alemania, y explica por qué el costo recae en los entrantes marginales, los perceptores secundarios de ingreso, los jóvenes y las mujeres, y no en los trabajadores protegidos que ya tenían un puesto. El aumento simultáneo de las horas es coherente con el paso hacia el trabajo independiente, donde las horas están menos restringidas. Un matiz acota el destino del margen desplazado: en el diseño regional, la tasa incondicional de trabajo independiente no aumenta con la exposición (Cuadro~\ref{tab:uncond} del apéndice), de modo que entre departamentos el margen absoluto robusto es la destrucción de empleo formal, y el aumento de la participación del trabajo independiente en el empleo refleja en parte la reducción de la base de empleo; el desplazamiento absoluto hacia el trabajo por cuenta propia aparece en el contraste por educación, donde se compara directamente a los trabajadores desplazados con sus sustitutos más cercanos.

¿Por qué la reasignación no estuvo acompañada de una respuesta salarial agregada? No porque el piso sea letra muerta: la evidencia de primera etapa muestra que, dentro del sector formal, el pico en el mínimo antiguo migró casi por completo al nuevo, de modo que los empleadores cubiertos cumplieron. La explicación está en la aritmética de la cobertura en un mercado laboral con dos sectores \citep{ulyssea_2018}. Solo el 13 por ciento de los trabajadores poco calificados ocupados es formal, así que incluso el cumplimiento pleno dentro del segmento formal mueve muy poco el salario promedio del grupo; mientras tanto, la remuneración en el segmento informal se alejó \emph{aún más} por debajo del piso, y ambas respuestas se compensan en el agregado. La inflación agrava el cuadro: el aumento de 2022, de diez por ciento en términos nominales, coincidió con una inflación de precios al consumidor de alrededor de ocho por ciento, de modo que el valor real del piso apenas subió y luego se erosionó. Un piso que solo se respeta donde los costos laborales son fiscalizables no puede elevar la distribución salarial del grupo, pero sí puede influir en los arreglos de empleo que eligen empresas y trabajadores, y según mis estimaciones lo hizo. Es la misma lógica de reasignación que \citet{dustmann_2022} documentan entre establecimientos en Alemania, trasladada al margen que importa en una economía en desarrollo: la frontera entre el trabajo formal y el informal \citep{jales_2018}.

Dos salvedades acotan esta lectura. La reasignación, aunque robusta a la especificación y a la inferencia basada en el diseño, es sensible a violaciones moderadas de las tendencias paralelas en el contraste por educación \citep{rambachan_roth_2023}, y no reaparece en torno a la reforma de 2025, donde el mismo diseño y las mismas magnitudes arrojan efectos nulos en todos los márgenes y una pequeña estimación de signo opuesto para el trabajo independiente. Interpreto el contraste entre ambos episodios como dependencia del estado, y la Sección~\ref{sec:reconciliation} examina directamente las dos variables de estado candidatas: la evidencia favorece la posición del piso en la distribución de la remuneración formal (el aumento de 2022 barrió una franja que concentraba la cuarta parte de la remuneración formal de baja educación, mientras que el de 2025 barrió una más delgada) y rechaza la versión basada en la rotación del mercado laboral, pues la creación de vínculos formales era de hecho mayor en vísperas de 2025. La evidencia del margen de entrada de la Sección~\ref{sec:transitions} identifica el margen por el que debe operar cualquier dependencia del estado de este tipo: la creación de nuevos vínculos formales, cuya rentabilidad gobierna el piso. Sin embargo, no puedo descartar del todo la lectura alternativa, según la cual una dinámica residual de recuperación correlacionada con la educación contribuye a las estimaciones de composición de 2022; lo que me inclina hacia una interpretación causal es la evidencia de dosis-respuesta y de exposición continua, pues una recuperación diferenciada por educación no tiene por qué seguir el índice de Kaitz regional previo a la reforma.

\subsection{Conciliación con la literatura previa}\label{sec:reconciliation}

Mis resultados encajan bien con parte de la evidencia previa y entran en tensión con el resto, y ese patrón es informativo en sí mismo. Hacen eco del trabajo peruano anterior de \citet{jaramillo_lopez_2006}, quienes encontraron que la reforma de 2003 fue en gran medida inocua, con desplazamiento hacia la informalidad, y de \citet{cespedes_2005}, quien encontró efectos concentrados entre los jóvenes. Son coherentes con el hallazgo internacional general de que los aumentos del salario mínimo en este rango producen efectos pequeños sobre el empleo \citep{card_krueger_1995, cengiz_2019}.

La tensión es con el estudio peruano más reciente. \citet{jimenez_2023} encuentra que la reforma de 2016 elevó los salarios formales entre cuatro y cinco por ciento y aumentó la formalidad, un patrón de efecto faro que no reproduzco para 2022. Tres diferencias explican plausiblemente la divergencia. El aumento de 2016 ocurrió en un periodo de baja inflación, por lo que su incidencia real fue mayor y más duradera que la del aumento de 2022, erosionado por la inflación. La identificación difiere: \citet{jimenez_2023} usa una dosis a nivel departamental definida sobre la distribución salarial formal, mientras que yo contrasto a trabajadores de baja y alta educación a nivel individual, un diseño que da más peso al amplio segmento informal de los poco calificados. Y el contexto macroeconómico difiere: crecimiento estable en 2016 frente a una recuperación pospandemia en 2022. Una observación ayuda a discriminar entre estas explicaciones: 2016, el episodio con efecto faro, fue también un año tranquilo, al igual que 2025, en el que no encuentro efecto alguno, de modo que el estado macroeconómico por sí solo no puede explicarlo todo. Lo que separa a 2016 de 2025 es el \emph{nivel} del piso cuando llegó el aumento: en 2016 el mínimo estaba bastante por debajo del salario mediano y un aumento real podía plausiblemente propagarse como norma, mientras que hacia 2025 el piso ya superaba la mediana nacional, de modo que un nuevo aumento partía de un nivel que la práctica ya había dejado de respetar para el trabajador marginal. Bajo esta lectura, los tres episodios se ordenan según dos variables de estado: cuán alto se ubica ya el piso en la distribución salarial y cuánta reasignación está atravesando el mercado laboral. Un piso moderado que sube en un mercado estable puede elevar las normas (2016); un piso alto que sube en medio de la reasignación pospandemia mueve el margen de composición (2022); y un piso alto que sube en un mercado asentado no mueve nada (2025).

Como una historia con dos variables de estado libres puede racionalizar casi cualquier cosa, examino ambas directamente y dejo que los datos descarten una. La variable de posición del piso se comporta como la historia exige: en vísperas de la reforma de 2022, el 24 por ciento de los asalariados formales de baja educación ganaba dentro de la franja que el nuevo piso estaba por barrer y el 16 por ciento se acumulaba en el piso antiguo, mientras que en vísperas de la reforma de 2025 esas proporciones habían bajado a 17 y 14 por ciento. La variable de reasignación no: la proporción de nuevas contrataciones entre los trabajadores formales de baja educación era \emph{mayor} en 2024 (41 por ciento) que antes de la reforma de 2022 (32 por ciento), de modo que el resultado nulo de 2025 no puede atribuirse a una escasez de nuevos vínculos formales. Una prueba de dosis a nivel departamental apunta en la misma dirección: interactuar el contraste por educación con la creación de vínculos formales previa a la reforma no produce ningún gradiente en ninguno de los dos episodios, con o sin el control de Kaitz.\footnote{Los momentos de las variables de estado y las estimaciones de dosis se incluyen en el paquete de replicación (\texttt{state\_premise.csv}, \texttt{state\_dose.csv}).} Por ello descarto la versión del argumento basada en la rotación y apoyo la interpretación en la posición del piso: el margen de entrada se cierra cuando el nuevo piso barre una franja densa de la remuneración formal, como en 2022, y le queda poco por cerrar cuando el piso ya superó el rango que la práctica de remuneración formal respeta, como en 2025. El contraste en la franja barrida es real pero modesto, por lo que presento esta interpretación como la más coherente con la evidencia y no como un hecho establecido.

\subsection{Validez externa y limitaciones}

Cinco limitaciones acotan la interpretación. Primero, la identificación descansa en el supuesto de que, sin la reforma, los trabajadores de baja y alta educación habrían seguido trayectorias paralelas. Pongo a prueba este supuesto donde es posible y acoto las consecuencias de que no se cumpla; las tendencias previas de la formalidad son planas, pero la prueba de tendencias previas salariales rechaza debido a un único trimestre de inicios de 2021, y la serie de empleo muestra una clara tendencia previa de recuperación, razón por la cual formulo la conclusión salarial como ``ninguna ganancia detectable'' y no como un nulo preciso, y trato el resultado sobre el nivel de empleo como no concluyente. Segundo, el diseño identifica efectos sobre los trabajadores poco calificados en relación con los calificados. El grupo de comparación no está del todo libre de exposición (el 40 por ciento de los asalariados calificados ganaba por debajo del nuevo piso antes de la reforma), de modo que las estimaciones son magnitudes relativas que, en primer orden, subestiman el efecto sobre los plenamente expuestos; el diseño de exposición continua aborda en parte este problema al prescindir del grupo de comparación. Tercero, el diseño principal usa cortes transversales repetidos; el análisis con el Panel ENAHO de la Sección~\ref{sec:transitions} aborda directamente la distinción entre flujos y stocks, pero es anual y no trimestral, su referencia previa a la reforma incluye transiciones desde el año pandémico 2020 y sus muestras de base formal son modestas, por lo que lo trato como una corroboración del mecanismo y no como fuente de magnitudes precisas. Cuarto, la informalidad en 2024 se mide con un indicador reconstruido y no con el oficial; dentro de la ventana de validación de 2025 el indicador se reconstruye de manera consistente en todos los trimestres, pero el empalme de 2024 añade error de medición al tramo final del estudio de eventos de 2022. Finalmente, la ventana abarca como máximo diez trimestres limpios posteriores a la reforma, por lo que me refiero a efectos de corto y mediano plazo y no al ajuste de largo plazo del capital ni a la entrada y salida de empresas.

\subsection{Implicancias de política}

Los resultados dejan un mensaje sobrio para la política de salario mínimo en una economía de alta informalidad. Un piso que no se hace cumplir para la mayoría de sus beneficiarios previstos hace poco por elevar su remuneración, y donde sí es vinculante, en el sector formal, puede empujar a algunos trabajadores hacia arreglos menos protegidos. Elevar el salario mínimo es, por tanto, un instrumento imperfecto para aumentar los ingresos de los poco calificados en el Perú, conclusión que hace eco de \citet{jaramillo_lopez_2006}. Esto no significa que el salario mínimo sea inocuo ni inútil; mis estimaciones son relativas y mi ventana es corta. Pero sí implica que la formalización y la fiscalización son condiciones previas para que el piso llegue a los trabajadores que busca proteger. Sin ellas, el margen de ajuste es la composición del empleo y no el nivel de la remuneración.
